% =====================================================================
%  A Semantic Model of Genetic Evidence
%
%  Extended version (A): self-contained, standard article class, for arXiv.
%  Migrated off the CEUR ceurart class (see manuscript-plan-jamia.md, Step 0).
% =====================================================================

\documentclass[11pt]{article}

\usepackage[margin=1in]{geometry}

% natbib + authblk are loaded here, NOT in the shared preamble: preamble.tex is
% also \input by supplement.tex, which already loads natbib, so loading it there
% too would double-load it. natbib must precede hyperref (loaded in the
% preamble), so it is loaded before % ---------- core packages ----------
\usepackage[T1]{fontenc}
\usepackage{microtype}        % better typography; load early
\usepackage{graphicx}
\usepackage{booktabs}         % nicer tables (\toprule, \midrule, \bottomrule)
\usepackage{tabularx}         % auto-width columns
\usepackage{array}
\usepackage{longtable}        % multi-page tables (UMLS/OMOP crosswalk section)
\usepackage{listings}         % code/YAML listings — chosen over minted for portability
\usepackage{xcolor}
\usepackage{tikz}
\usetikzlibrary{positioning, arrows.meta, calc}
\usepackage{amsmath,amssymb}
\usepackage[section]{placeins} % flushes floats at every \section boundary
\usepackage{hyperref}         % load near last
\usepackage{cleveref}         % load after hyperref; \cref{} auto-adds "Figure"/"Table"/etc.

% ---------- breakable file paths ----------
% \fpath{...} typesets a file path in typewriter that may wrap at path
% separators, so long paths (e.g. schema/genetic_evidence.shacl.ttl) do
% not run past the right margin. Verbatim-style: do NOT escape "_".
\DeclareUrlCommand\fpath{\urlstyle{tt}}
\Urlmuskip=0mu plus 1mu\relax      % allow a little stretch around break points
\makeatletter
\g@addto@macro\UrlBreaks{\do\.\do\-}  % break at . and - as well as the default /
\makeatother

% ---------- listings setup for YAML annotations ----------
% Define a lightweight YAML "language" for listings.
\definecolor{yamlKey}{RGB}{0,90,156}
\definecolor{yamlStr}{RGB}{163,21,21}
\definecolor{yamlComment}{RGB}{0,128,0}
\definecolor{listingBg}{RGB}{248,248,248}

\lstdefinelanguage{yaml}{
  keywords={true,false,null},
  keywordstyle=\color{yamlKey}\bfseries,
  sensitive=true,
  comment=[l]{\#},
  commentstyle=\color{yamlComment}\itshape,
  stringstyle=\color{yamlStr},
  morestring=[b]',
  morestring=[b]",
  moredelim=[l][\color{yamlKey}]{-\ },
}

\lstset{
  language=yaml,
  inputencoding=utf8,
  extendedchars=true,
  literate={—}{{---}}1 {–}{{--}}1 {§}{{\S}}1 {β}{{$\beta$}}1 {→}{{$\rightarrow$}}1 {₂}{{\textsubscript{2}}}1,
  basicstyle=\ttfamily\small,
  backgroundcolor=\color{listingBg},
  frame=single,
  framerule=0pt,
  rulecolor=\color{listingBg},
  xleftmargin=1em,
  xrightmargin=1em,
  aboveskip=0.8em,
  belowskip=0.8em,
  showstringspaces=false,
  columns=fullflexible,
  breaklines=true,
  breakatwhitespace=true,
  captionpos=b,
}

% ---------- typography commands for our model's core terms ----------
% Using \textsc for the abstract classes so they stand out as technical terms
% without being as loud as monospace.
\newcommand{\ScientificEvidence}{\textsc{ScientificEvidence}}
\newcommand{\EvidenceVariable}{\textsc{EvidenceVariable}}
\newcommand{\EvidenceAssertion}{\textsc{EvidenceAssertion}}
\newcommand{\GeneticEvidence}{\textsc{GeneticEvidence}}
\newcommand{\GeneticEvidenceVariable}{\textsc{GeneticEvidenceVariable}}
\newcommand{\GeneticEvidenceAssertion}{\textsc{GeneticEvidenceAssertion}}

% For categorical values of dimensions (e.g., STATISTICAL GENETICS), we use
% small caps so they are visually distinct from prose but not as heavy as code.
\newcommand{\dimval}[1]{\textsc{\lowercase{#1}}}
.
% Cross-document references into the supplement: xr imports a labels-only aux
% (target/supp-labels.aux, extracted by compile_paper.sh, which builds the
% supplement first). \ref* (hyperref's non-linking form) renders the number
% without attempting an intra-document hyperlink.
\usepackage{xr}
\externaldocument[supp:]{target/supp-labels}
\usepackage[numbers,sort&compress]{natbib}
\usepackage{authblk}

% ---------- shared preamble (factored out; also \input by supplement.tex) ----------
% ---------- core packages ----------
\usepackage[T1]{fontenc}
\usepackage{microtype}        % better typography; load early
\usepackage{graphicx}
\usepackage{booktabs}         % nicer tables (\toprule, \midrule, \bottomrule)
\usepackage{tabularx}         % auto-width columns
\usepackage{array}
\usepackage{longtable}        % multi-page tables (UMLS/OMOP crosswalk section)
\usepackage{listings}         % code/YAML listings — chosen over minted for portability
\usepackage{xcolor}
\usepackage{tikz}
\usetikzlibrary{positioning, arrows.meta, calc}
\usepackage{amsmath,amssymb}
\usepackage[section]{placeins} % flushes floats at every \section boundary
\usepackage{hyperref}         % load near last
\usepackage{cleveref}         % load after hyperref; \cref{} auto-adds "Figure"/"Table"/etc.

% ---------- breakable file paths ----------
% \fpath{...} typesets a file path in typewriter that may wrap at path
% separators, so long paths (e.g. schema/genetic_evidence.shacl.ttl) do
% not run past the right margin. Verbatim-style: do NOT escape "_".
\DeclareUrlCommand\fpath{\urlstyle{tt}}
\Urlmuskip=0mu plus 1mu\relax      % allow a little stretch around break points
\makeatletter
\g@addto@macro\UrlBreaks{\do\.\do\-}  % break at . and - as well as the default /
\makeatother

% ---------- listings setup for YAML annotations ----------
% Define a lightweight YAML "language" for listings.
\definecolor{yamlKey}{RGB}{0,90,156}
\definecolor{yamlStr}{RGB}{163,21,21}
\definecolor{yamlComment}{RGB}{0,128,0}
\definecolor{listingBg}{RGB}{248,248,248}

\lstdefinelanguage{yaml}{
  keywords={true,false,null},
  keywordstyle=\color{yamlKey}\bfseries,
  sensitive=true,
  comment=[l]{\#},
  commentstyle=\color{yamlComment}\itshape,
  stringstyle=\color{yamlStr},
  morestring=[b]',
  morestring=[b]",
  moredelim=[l][\color{yamlKey}]{-\ },
}

\lstset{
  language=yaml,
  inputencoding=utf8,
  extendedchars=true,
  literate={—}{{---}}1 {–}{{--}}1 {§}{{\S}}1 {β}{{$\beta$}}1 {→}{{$\rightarrow$}}1 {₂}{{\textsubscript{2}}}1,
  basicstyle=\ttfamily\small,
  backgroundcolor=\color{listingBg},
  frame=single,
  framerule=0pt,
  rulecolor=\color{listingBg},
  xleftmargin=1em,
  xrightmargin=1em,
  aboveskip=0.8em,
  belowskip=0.8em,
  showstringspaces=false,
  columns=fullflexible,
  breaklines=true,
  breakatwhitespace=true,
  captionpos=b,
}

% ---------- typography commands for our model's core terms ----------
% Using \textsc for the abstract classes so they stand out as technical terms
% without being as loud as monospace.
\newcommand{\ScientificEvidence}{\textsc{ScientificEvidence}}
\newcommand{\EvidenceVariable}{\textsc{EvidenceVariable}}
\newcommand{\EvidenceAssertion}{\textsc{EvidenceAssertion}}
\newcommand{\GeneticEvidence}{\textsc{GeneticEvidence}}
\newcommand{\GeneticEvidenceVariable}{\textsc{GeneticEvidenceVariable}}
\newcommand{\GeneticEvidenceAssertion}{\textsc{GeneticEvidenceAssertion}}

% For categorical values of dimensions (e.g., STATISTICAL GENETICS), we use
% small caps so they are visually distinct from prose but not as heavy as code.
\newcommand{\dimval}[1]{\textsc{\lowercase{#1}}}


% ---------- front matter ----------
\begin{document}

    \title{A Semantic Model of Genetic Evidence: A Step Toward Bridging the Basic-Science--Clinic Gap}

% Authors (authblk). ORCIDs retained as comments; add an \orcid package for a
% camera-ready if desired.
    \author[1,2,3]{Michael Bouzinier\thanks{Correspondence: \texttt{michael\_bouzinier@harvard.edu}}}
    \author[3,4]{Dmitry Etin}
% ORCIDs: Bouzinier 0000-0002-3161-5601; Etin 0000-0003-1068-2781.
    \affil[1]{Harvard University, Cambridge, MA, USA}
    \affil[2]{IDEXX Laboratories, Westbrook, ME, USA}
    \affil[3]{Forome Association, Newton, MA, USA}
    \affil[4]{Deggendorf Institute of Technology, Germany}
    \date{}

    \maketitle

    \begin{abstract}
        Scientific and clinical decision-making depends on evidence from
        the primary literature, but existing standards for representing
        that evidence (FHIR Evidence, ECO, SEPIO, and the GA4GH Genomic
        Knowledge Standards) are oriented toward clinical-trial workflows,
        evidence codes, or single-variant assertions, and do not capture
        the fine-grained, domain-specific structure of claims in basic and
        pre-clinical research. We introduce a semantic model for scientific
        evidence with three core classes, specialize it for genetics, align it structurally to FHIR Evidence with a
        SEPIO-anchored credibility decomposition, and attach a compact
        dimensional vocabulary whose conditional-activation rules are
        validated by a SHACL schema for the implemented constraints. Using clinical variant
        interpretation as the driving use case, we evaluate the model
        through a human--AI annotation pilot over six genetics papers,
        yielding 28 evidence items and 95 source-anchored assertions, with
        a workflow that keeps curator-authored ground truth distinct from
        AI-drafted annotations. Treating the pilot as a feasibility study
        rather than a benchmark, we argue that the model is a useful
        increment toward trustworthy, AI-ready infrastructure for variant
        interpretation: a reference data model and validation schema for
        representing genetic evidence.
    \end{abstract}

    \medskip
    \noindent\textbf{Keywords:} genetic evidence; variant interpretation;
    ontology; FHIR; ECO; SEPIO; GA4GH; AI-ready data.

% =====================================================================
% §1 — Introduction (see paper/sections/01_intro.tex)
    % =====================================================================
%  §1 - Introduction
%
% =====================================================================

\section{Introduction}
\label{sec:intro}

Scientific and clinical decision-making depends on evidence from
heterogeneous sources, especially the primary literature. Today,
this evidence is largely consumed by humans and encoded indirectly
into annotation databases and local rules.

Existing resources such as HL7 Fast Healthcare Interoperability
Resources (FHIR) Evidence and EvidenceVariable~\cite{fhir_evidence},
the Evidence and Conclusion Ontology (ECO)~\cite{eco2019}, the
Scientific Evidence and Provenance Information Ontology
(SEPIO)~\cite{sepio}, and the Global Alliance for Genomics and
Health (GA4GH) Genomic Knowledge Standards~\cite{ga4gh_va}
advance interoperability in clinical research. However, they do not
comprehensively represent the evidence types and workflows of the
basic sciences. They focus mainly on clinical trials, high-level
assertions, and variant representations, with limited support for the
fine-grained, domain-specific evidence from basic and pre-clinical
research needed for automated reasoning and AI-ready infrastructure.

Each scientific domain demands a tailored approach for structuring
and communicating evidence. For example, the genetic evidence
discussed in this paper is generated in diverse contexts, from
experimental model systems to large-scale population studies, that
are not fully supported by current standards.
To properly use evidence derived from
scientific literature in clinical workflows, two principal challenges
require standardisation: \emph{Relevance and Specificity}, that is,
how precisely a genetic finding maps to an individual patient,
disease context, or biological hypothesis; and \emph{Confidence},
that is, whether a given piece of evidence can be robustly and
safely used, ensuring defensible compliance, interpretability, and
reproducibility.

A well-designed schema for genetic evidence helps make literature
review a \emph{semantic-parsing} problem: mapping natural-language
claims to a shared, source-anchored representation. Large language
models are increasingly capable of this kind of parsing, but
dependable automation still requires a carefully structured target
language. The present work offers one such increment. Building on decades of work by
the biomedical-ontology community, it takes a small step toward literature
review that is materially assisted by automated parsing.

\paragraph{Contributions.} This paper makes the following contributions:
\begin{itemize}
    \item A domain-analysis-driven semantic model for genetic evidence, with
    core classes (\ScientificEvidence, \EvidenceVariable,
    \EvidenceAssertion), genetic specialisations, a compact dimensional
    vocabulary, and a validation schema in the W3C Shapes Constraint
    Language (SHACL)~\cite{w3c-shacl}.
    \item An application of the model to six genetic-evidence papers spanning
        human, population, model-organism, and polygenic-score genetics,
        supported by extraction tooling.
    \footnote{\url{https://github.com/ForomePlatform/genetic-evidence-model}}
    \item A human--AI collaborative annotation workflow, packaged as executable
        AI skills for annotation and review that apply the model under curator
        supervision.
    \item A structural alignment of the core classes with HL7 FHIR Evidence,
        including a credibility representation that decomposes certainty into
        SEPIO-anchored facets. The dimensions are designed to align with ECO,
        SEPIO, IAO, and OBI; Supplementary Note~SN5 provides a partial
        term-level mapping.
    \item A curated term-level UMLS crosswalk of the dimensional vocabulary, in
        which every value is mapped with a typed relation or recorded as an
        argued gap. The crosswalk was produced with the open-source,
        model-agnostic \texttt{GEM Mapping Studio}, a standalone curation tool
        distributed as a Python package that supports the full workflow, from
        defining dimension axes in an empty workspace to adjudicating
        value-level mappings. The workflow follows a reproducible protocol with
        pre-registered search sweeps, structured rejection rationales, and a
        public decision log in the repository.
\end{itemize}

The remainder of the paper reviews related standards
(\Cref{sec:background}), defines the semantic model
(\Cref{sec:model}), and illustrates it with two worked examples
(\Cref{sec:examples}). It then describes the annotation workflow and
pilot statistics (\Cref{sec:workflow}), discusses findings, candidate
extensions, limitations, and future work (\Cref{sec:discussion}), lists
the released artefacts (\Cref{sec:availability}), and concludes
(\Cref{sec:conclusion}).


% =====================================================================

% =====================================================================
% §2 — Background and Related Work (see paper/sections/02_background.tex)
    \section{Background and Relation to Existing Standards}
\label{sec:background}

The representation of scientific and clinical evidence has been developed
across clinical informatics, biomedical ontology, and genomic knowledge
representation. Seven families of standards bear directly on our work: HL7 FHIR
Evidence~\cite{fhir_evidence}; the Evidence and Conclusion Ontology
(ECO)~\cite{eco2019} and the Scientific Evidence and Provenance Information
Ontology (SEPIO)~\cite{sepio}; the GA4GH Variant Annotation (VA) and Variation
Representation (VRS) specifications~\cite{ga4gh_va,ga4gh_vrs}; the ACMG/AMP
variant-interpretation guidelines~\cite{acmg2015}; the upper ontologies IAO
and OBI~\cite{iao2015,obi2016}; the PROV-O provenance model~\cite{prov_o};
and the ClinGen, MONDO, and HPO
curation and terminology resources~\cite{clingen2018,mondo,hpo}. Each captures
part of what a basic-science genetic publication reports, but none captures the
whole: a single conditional-knockout study may report molecular, cellular, and
clinical phenotypes at once, and a polygenic-score study may aggregate millions
of variants into one score. A per-family review, with the specific coverage
gaps that motivate our model, is given in Supplementary Note~\ref*{supp:sec:supp-background}. Of these
families, FHIR Evidence is the closest structural relative of our model, and we
position the model relative to it here.

\subsection{Relation to FHIR Evidence}
\label{sec:background-fhir}

FHIR is an HL7 information model rather than an OBO Foundry ontology. A
genetic-evidence model could, in principle, be built directly on the OBO
side, but doing so first would be premature. ECO and SEPIO provide,
respectively, an evidence-type vocabulary and a model of assertion
provenance and qualifiers. Neither provides the variable-and-certainty
content model for an evidence unit that FHIR Evidence provides and that
basic-science evidence needs to support translation from bench to bedside.

Aligning a genetics-only model directly to OBO before such a general content
model exists would yield a narrow resource. It would neither compose cleanly
with other basic-science domains nor reuse OBO's orthogonal vocabularies,
creating exactly the redundancy that OBO's emphasis on orthogonality and
reuse is intended to prevent.

We therefore take the structural content model from FHIR, generalize it from
clinical trials to basic science, and design the dimensions to align with
ECO, SEPIO, IAO, and OBI. Supplementary Note~\ref*{supp:sec:supp-crosswalk} provides a term-level
crosswalk to these vocabularies, including verified identifiers and the
remaining gaps. Genetics is the first specialization; extending the model to
other basic-science domains and completing the term-level alignment are
future work.

The HL7 FHIR Evidence, EvidenceVariable, and EvidenceReport
resources~\cite{fhir_evidence} were designed with systematic-review and
guideline-development workflows in mind, where evidence concerns an
intervention tested in a population. FHIR Evidence therefore has a different native
framing and does not by itself represent the evidence produced by basic and
pre-clinical research. Our three core classes address this gap.

The core classes are general across the basic sciences rather than specific
to genetics. They retain FHIR's separation between the evidence variable and
the certainty of a finding, incorporate FHIR's separate \texttt{statistic}
into \EvidenceVariable{} as the measured value, and replace the
clinical-trial content frame of population, intervention, comparator, and
outcome with a genetics-specific dimensional frame. Genetics is the
specialization developed in this paper (\cref{sec:model}):
\ScientificEvidence{}, \EvidenceVariable{}, and \EvidenceAssertion{}
specialize to \GeneticEvidence{}, \GeneticEvidenceVariable{}, and
\GeneticEvidenceAssertion{}, respectively. \Cref{tab:fhir-alignment} gives
the class-level correspondence. The field-level alignment, including the
certainty subcomponents discussed in \cref{sec:model-dimensions}, is given in
Supplementary Table~\ref*{supp:tab:fhir-fields}.
\footnote{We distinguish two notions that the literature often conflates.
The \texttt{assertion} property of a \ScientificEvidence{} item records the
object-level claim made by a paper. It is what the pilot records, and its
closest analogues are the GA4GH VA \texttt{Statement} and the SEPIO
\texttt{Assertion}, not the FHIR \texttt{Evidence.assertion} element, which
is a human-readable summary. \EvidenceAssertion{} is instead a
meta-predicate over decision rules that would use such
claims~\cite{our_preprint}. It has no direct analogue in FHIR, VA, or SEPIO and
is not instantiated in this paper. The property and the class share a name;
a future schema revision will resolve this collision.}

\begin{table}[tbp]
    \centering
    \caption{Class-level correspondence between the core model and FHIR Evidence
    R5. The alignment is architectural: the genetics-specific dimensional frame
    replaces FHIR's clinical-trial frame. Field-level detail is provided in
    Supplementary Table~\ref*{supp:tab:fhir-fields}.}
    \label{tab:fhir-alignment}
    \small
    \begin{tabularx}{\linewidth}{@{}llX@{}}
\toprule
This work & FHIR Evidence R5 & Relationship \\
\midrule
\ScientificEvidence
& \texttt{Evidence}
& Structural sibling; GEM replaces the PICO frame with a
  genetics-specific dimensional frame \\
\EvidenceVariable
& \texttt{EvidenceVariable} + \texttt{Evidence.statistic}
& GEM combines the variable definition and measured value in one class \\
\EvidenceAssertion
& no equivalent
& Meta-predicate over decision rules; part of the reasoning layer
  described in~\cite{our_preprint}, which is outside the scope of this paper \\
Credibility
& \texttt{Evidence.certainty}
& Structural parallel; both provide an overall rating with separable,
  genetics-specific subcomponents \\
\bottomrule
    \end{tabularx}
\end{table}


% =====================================================================

% =====================================================================
% =====================================================================
% §3 — Model (see paper/sections/03_model.tex)
    % =====================================================================
%  §3 — A Semantic Model of Genetic Evidence
%
%  Core classes, dimensional vocabulary, conditional activation
%  mechanism. Written to be self-contained; subsequent sections
%  reference back to the table and the formal snippets here.
% =====================================================================

\section{A Semantic Model of Genetic Evidence}
\label{sec:model}

The model contains three parts: a small set of \emph{core classes}
(\cref{sec:model-classes}) that define what a unit of evidence is and
how evidence items compose; a \emph{dimensional vocabulary}
(\cref{sec:model-dimensions}) that describes each evidence item along
roughly twenty axes; and a \emph{conditional-activation} mechanism
(\cref{sec:model-conditional}) that keeps the vocabulary compact by
requiring dimensions only when they are semantically appropriate.

\subsection{Core Classes}
\label{sec:model-classes}

Three abstract classes form the core of the model:

\begin{description}
    \item[\ScientificEvidence] A unit of evidence, typically linked to a
          publication or dataset and described by a set of dimensions (see
          \cref{tab:dimensions-core,tab:dimensions-cond}). It carries an
          \texttt{assertion} property containing the object-level,
          source-anchored claims made by the source, for example that a
          variant is associated with a phenotype. These are the claims
          recorded in the pilot (\cref{sec:workflow}). A single paper may
          give rise to multiple \ScientificEvidence{} items when it reports
          distinct evidential claims, such as a case--control association and
          an \emph{in vivo} functional assay of the same variant.
    \item[\EvidenceVariable] A named quantity extracted from one or more
        \ScientificEvidence{} items, such as the gnomAD allele frequency of
        a variant, the LOD score in a family, or the odds ratio reported by
        a GWAS. Evidence variables are the terms from which predicates over
        Evidence are constructed. They are also described by dimensions,
        because the same quantity, such as an odds ratio, may carry different
        epistemic weight depending on how it was obtained.
    \item[\EvidenceAssertion] A meta-predicate, that is, a predicate about a
        predicate. It determines whether \EvidenceVariable{}s are used
        appropriately in a predicate over Evidence, namely in a decision
        rule. This class is the model's point of contact with the reasoning
        layer described in~\cite{our_preprint}, where such meta-predicates
        assert which evidence a rule may use. This paper constructs no
        decision rules and therefore instantiates no \EvidenceAssertion{}.
        The class remains part of the model, but the interpretability and
        AI-readiness of the reasoning layer are the primary focus
        of~\cite{our_preprint}, not of the present pilot.
\end{description}

The \texttt{assertion} property and the \EvidenceAssertion{} class share
a name but denote different things. The property records the object-level
claims captured in the pilot, whereas the class is a meta-predicate over
any decision rule that would use those claims. The property name is
retained for compatibility with the released annotations and will be
renamed, for example to \texttt{selector}, in a future schema revision.

Specialising these classes for the genetic domain yields
\GeneticEvidence, \GeneticEvidenceVariable, and
\GeneticEvidenceAssertion. These specialisations inherit the provenance
requirement and add the dimensional annotations described next.

\subsection{Dimensions}
\label{sec:model-dimensions}

Each \GeneticEvidence{} instance is described using a dimensional
vocabulary summarised in \cref{tab:dimensions-core}, which lists the
always-required dimensions, and \cref{tab:dimensions-cond}, which gives
representative conditional dimensions.\footnote{The full enumeration of
conditional dimensions, including those omitted from
\cref{tab:dimensions-cond} for space, is maintained in the project
repository at
\url{https://github.com/ForomePlatform/genetic-evidence-model/blob/master/schema/dimensions.md}.}

Six dimensions are \emph{always required}: Knowledge Domain, Method,
Target Type, Resolution, Credibility, and Phenotype Scale. No
well-formed genetic-evidence item is fully specified without them.
Together, they tell a downstream reader what kind of evidence the item
is, how it was obtained, what it concerns, and how credible it is in
context.

The dimensions differ in internal structure, and each is modeled
according to its own semantics rather than forced into a single
hierarchy. Knowledge Domain values are sibling framings rather than a
subsumption hierarchy. A single item may carry several values at once,
as when a functional study is interpreted in a clinical-genetics
context; the multi-valued cardinality captures this directly.

Phenotype Scale is an ordinal scale, from molecular to clinical, rather
than a class hierarchy. Modeling it through subsumption would therefore
misrepresent the dimension. Method, by contrast, is genuinely
hierarchical, and its \texttt{is\_a} relations are encoded in the SHACL
distribution. The chosen structure is therefore part of the definition
of each dimension.

Phenotype Scale and Variant Ascertainment were added during the annotation
work (\cref{sec:workflow}). Phenotype Scale distinguishes molecular
phenotypes, such as binding or activity, from cellular, histological,
organismal, and clinical phenotypes. It is independent of the scale of the
genetic target.

Variant Ascertainment records how a variant entered a study. It is meaningful
only for variant-level targets and is therefore a \emph{conditional}
dimension, required when Target Type is \dimval{Variant}. Both dimensions
are revisited in \cref{sec:discussion}.

Credibility records the methodological soundness of an asserted finding, and
nothing else. It is a structured decomposition rather than an open key-value
field: one overall ordinal rating together with separately rated facets for
statistical power or sample size, independent replication, multiple-testing
control, ancestry or population-stratification control, and ascertainment.
The facets remain separate rather than being collapsed into a single score
because clinical and research users may weigh them differently.

This separation is inspired by GRADE's general practice of assessing the
factors that affect certainty separately~\cite{grade2008}. GEM does not
adopt the GRADE framework or its rating rules, which were developed for
bodies of clinical evidence, particularly randomized and observational
studies. The GEM facets instead support credibility judgments about
individual evidence items from basic and pre-clinical research.

The overall rating is anchored on SEPIO confidence, defined as the degree of
belief that an asserted proposition is true~\cite{sepio}, and is implemented
as a genetics-specific value set bound to that attribute, as prescribed by
SEPIO profiles (Supplementary Table~\ref*{supp:tab:credibility-facets}).

Two constructs are excluded from credibility by design. First, the
direction or hedging of a finding belongs to the assertion itself and is
represented by the width of the value range it constrains. An honestly
reported null result from a rigorous study may therefore have high
credibility. Second, the extent to which evidence bears on a particular
question is applicability. It is captured by Relevance and Specificity and
computed for each question. Combining methodological soundness, assertion
content, and applicability in one field would conflate distinct constructs.

The remaining dimensions are \emph{conditional}: they are required, or
even meaningful, only when other dimensions take particular values. For
example, mode of inheritance is not meaningful for a polygenic score, and
knockout type is not meaningful for a human case--control study. Rather
than marking such dimensions as ``optional'' and relying on convention, the
model assigns explicit activation conditions to each, as described in
\cref{sec:model-conditional}. These conditions specify when a dimension
must be present; otherwise, it may be omitted. This makes the requirement
status transparent to both human curators and automated validators.

\begin{table}[tbp]
\centering
\caption{Always-required dimensions of the \GeneticEvidence{} model.
These apply to every evidence item regardless of context. Phenotype Scale
was added to this set during the annotation work
(\cref{sec:workflow}). KD abbreviates Knowledge Domain.}
\label{tab:dimensions-core}
\small
\begin{tabularx}{\linewidth}{@{}llllX@{}}
\toprule
Dimension              & Value type        & Cardinality & Objectivity & Example values \\
\midrule
Knowledge Domain       & categorical       & multiple & obj.       & \dimval{Human Genetics}, \dimval{Population Genetics}, \dimval{Model Organism}, \dimval{Gene Function} \\
Method                 & categorical       & multiple & obj.       & \dimval{Statistical Genetics}, \dimval{In Vivo}, \dimval{In Vitro}, \dimval{Clinical Evidence} \\
Target Type            & categorical       & single   & obj.       & \dimval{Gene}, \dimval{Variant}, \dimval{Interval}, \dimval{Transcript} \\
Resolution             & categorical       & single   & obj.       & \dimval{Window}, \dimval{Gene}, \dimval{Functional Element}, \dimval{Position}, \dimval{Variant} \\
Credibility            & rating ${+}$ facets & multiple & subj.      & overall rating ${+}$ \texttt{sample\_size}, \texttt{replication}, \texttt{multiple\_testing}, \texttt{stratification\_control}, \ldots \\
Phenotype Scale        & categorical       & single   & obj.       & \dimval{Molecular}, \dimval{Cellular}, \dimval{Histological}, \dimval{Organismal}, \dimval{Clinical} \\
\bottomrule
\end{tabularx}
\end{table}

\begin{table}[tbp]
\centering
\caption{Conditional dimensions of the \GeneticEvidence{} model.
Each is required only when its activation condition holds. KD, MT, and TT
abbreviate Knowledge Domain, Method, and Target Type, respectively;
$\supseteq$ denotes set membership for multi-valued dimensions.}
\label{tab:dimensions-cond}
\small
\begin{tabularx}{\linewidth}{@{}lllX@{}}
\toprule
Dimension              & Value type   & Cardinality & Activation condition \\
\midrule
Mode of Inheritance    & categorical  & single   & KD $\supseteq$ \dimval{Human Genetics} AND TT = \dimval{Gene} \\
Mendelian Segregation  & boolean      & single   & KD $\supseteq$ \dimval{Human Genetics} AND TT = \dimval{Gene} \\
Penetrance             & categorical  & single   & KD $\supseteq$ \dimval{Human Genetics} AND TT = \dimval{Variant} \\
Variant Ascertainment  & categorical  & multiple & TT = \dimval{Variant} \\
Measurement Target     & categorical  & multiple & KD $\supseteq$ \dimval{Gene Function} \\
Gene Relation          & categorical  & single   & KD $\supseteq$ \dimval{Gene Function} \\
Organism               & categorical  & multiple & MT $\supseteq$ \dimval{In Vivo} OR KD $\supseteq$ \dimval{Model Organism} \\
Knockout Type          & categorical  & single   & KD $\supseteq$ \dimval{Model Organism} \\
\bottomrule
\end{tabularx}
\end{table}

\subsection{Conditional Activation}
\label{sec:model-conditional}

Conditional activation is the least conventional part of the model, but it
keeps the vocabulary compact without sacrificing expressive power. Rather
than using a fixed schema in which every \GeneticEvidence{} item has every
dimension, with many dimensions marked as optional and left empty, the model
requires particular dimensions only when specified conditions over other
dimensions hold. This mirrors expert reasoning. A curator reading a paper
about an X-linked recessive pedigree asks about mode of inheritance because
the evidence makes that question meaningful, not because a global schema
requires it.

The activation conditions are encoded as SHACL shapes with SPARQL-based
constraints. This is the standard SHACL mechanism for expressing
requirements that cannot be captured by simple cardinality constraints on a
single property. Two representative shapes illustrate the pattern. For
brevity, we show only the \texttt{sh:sparql} body of each and omit the
outer NodeShape wrapper. The full shapes, together with the base
\texttt{GeneticEvidenceShape} and the simpler
\texttt{variant\_ascertainment} rule, are available in the project
repository with an annotated walkthrough.\footnote{\url{https://github.com/ForomePlatform/genetic-evidence-model/blob/master/schema/examples.md}}

Checking activation conditions is a closed-world question: whether an
annotation contains every dimension required by the evidence item. SHACL can
answer this question directly. OWL's open-world semantics cannot express
this form of completeness, which is one reason the constraints are encoded
in SHACL rather than as ontology axioms.

The first example encodes a conjunction across two classification axes.
\emph{Mode of Inheritance} is required when Knowledge Domain contains
\dimval{Human Genetics} \textbf{and} Target Type is \dimval{Gene}
(\cref{lst:shacl-moi}).

\begin{lstlisting}[language=yaml, caption={Compact form of the conditional-activation shape for \emph{Mode of Inheritance}. The outer NodeShape wrapper is omitted. The two triple patterns in \texttt{WHERE} form the conjunction; \texttt{FILTER NOT EXISTS} identifies instances that satisfy the condition but lack the required dimension.}, label={lst:shacl-moi}]
sh:sparql [ sh:select """
  SELECT $this WHERE {
    $this gem:knowledgeDomain gem:HUMAN_GENETICS .
    $this gem:targetType      gem:GENE .
    FILTER NOT EXISTS { $this gem:modeOfInheritance ?m }
  }""" ] .
\end{lstlisting}

The second example encodes a disjunction across two classification axes.
\emph{Organism} is required when Method contains \dimval{In Vivo}
\textbf{or} Knowledge Domain contains \dimval{Model Organism}
(\cref{lst:shacl-org}). The two disjuncts refer to different dimensions:
one is a Method value and the other is a Knowledge Domain value. This
illustrates that activation conditions can span the dimensional vocabulary
rather than being confined to a single axis.

\begin{lstlisting}[language=yaml, caption={Compact form of the conditional-activation shape for \emph{Organism}. The disjunction is expressed by a SPARQL \texttt{UNION} whose two branches refer to different dimensions.}, label={lst:shacl-org}]
sh:sparql [ sh:select """
  SELECT $this WHERE {
    { $this gem:method          gem:IN_VIVO . }
    UNION
    { $this gem:knowledgeDomain gem:MODEL_ORGANISM . }
    FILTER NOT EXISTS { $this gem:organism ?o }
  }""" ] .
\end{lstlisting}

In each case, the \texttt{SELECT} returns every \GeneticEvidence{}
instance that satisfies the activation condition but lacks the required
dimension. The validator reports these instances as constraint violations
with a localized diagnostic. Writing each conditional rule explicitly adds
some complexity, but it improves downstream clarity: a reader of an
annotation file can see why a particular dimension is required or absent,
and a validator can check conformance automatically.

\Cref{sec:workflow} returns to conditional activation in the context of
the annotation workflow.


% =====================================================================

% =====================================================================
% =====================================================================
% §4 — Worked Examples (see paper/sections/04_examples.tex)
    % =====================================================================
%  §4 — Worked Examples from the Literature
%
%  Six-paper corpus table, two featured examples (Duerr 2006, Inouye
%  2018) as short prose, with full annotations hosted as case reports
%  in the project repository.
%
%  Revised 2026-04-21: YAML listings moved to repository case-reports
%  to save space; per-example prose retained the methodological points
%  that the listings previously anchored.
% =====================================================================

\section{Illustrative Examples from the Literature}
\label{sec:examples}

To demonstrate the model's coverage and identify where it reaches its
current limits, we applied it to six published genetic-evidence papers that
span the major epistemic shapes encountered in literature-based variant
interpretation (Supplementary Table~\ref*{supp:tab:corpus}). The set covers
classical Mendelian human genetics (Gupta), population-scale association
(Duerr), case--control resequencing (Davis), model-organism and \emph{in
vitro} functional genetics (Jossin, Nelson, and the functional components of
Davis), and a genome-wide polygenic score (Inouye). Four of the six papers
were annotated manually from curator PDF highlights and callouts and form
the curator-authored reference set for this feasibility corpus. The remaining two were drafted by an AI assistant
and reviewed by a curator. All annotations and per-paper case reports are
available in the project
repository.\footnote{\url{https://github.com/ForomePlatform/genetic-evidence-model/tree/master/case-reports}}

Duerr et al.\ 2006 (\cref{sec:ex-duerr}) is the reference example of a
classical variant-level association study. Inouye et al.\ 2018
(\cref{sec:ex-inouye}) is a polygenic risk score study that exposes
limitations of the current variant-level model and motivates a cluster of
candidate extensions (\cref{sec:discussion}).

\subsection{A Genome-Wide Association Study Example}
\label{sec:ex-duerr}

Duerr et al.\ 2006 identified \textit{IL23R} as an inflammatory bowel
disease gene through a genome-wide association study of non-Jewish patients
with ileal Crohn's disease, followed by replication in a Jewish cohort and
family-based transmission testing. Under our decomposition, the paper yields
seven \GeneticEvidence{} items (Supplementary
Table~\ref*{supp:tab:corpus} and Case Report CR-DU): variant-level discovery
and replication, a family-based transmission analysis, conditional
fine-mapping, negative results at related pathway genes, and cited
functional and mouse-model support.

Six of the seven items fit the current schema without strain. The seventh, a
cited anti-p40 antibody trial, concerns drug-response evidence for which the
model has no fitting genetic Knowledge Domain, revealing a
translational-scope gap. Review of this annotation also identified three
candidate extensions that remain provisional under the two-papers rule: an
explicit direction-of-effect dimension for protective and risk-conferring
alleles, a construct for cross-item therapeutic synthesis, and an extension
of conditional activation from gene-level to variant-level Mendelian context.
Further detail is given in Supplementary
Note~\ref*{supp:sec:supp-ex-duerr}.

\subsection{A Polygenic Risk Score Example}
\label{sec:ex-inouye}

Inouye et al.\ 2018 developed and externally validated a meta-genomic risk
score (metaGRS) for coronary artery disease by combining three component
scores and evaluating performance in UK Biobank participants. Unlike Duerr's
variant-level association study, the evidential target is a derived,
whole-genome score and the primary measurements are epidemiological.

Our annotation decomposes the paper into four \GeneticEvidence{} items:
predictive and stratifying performance; independence from conventional risk
factors; persistence of stratification among medicated individuals; and the
headline finding that combining three component scores outperforms each
component (Case Report CR-IN). All four require deliberate placeholder
assignments in the current model. In particular, the score is assigned the
placeholder Target Type \dimval{Variant}, and Variant Ascertainment is marked
not applicable. These forced fits are recorded as six candidate extensions:
a score target type, whole-genome-aggregate resolution, a target-composition
axis, epidemiological measurement targets, a structured cohort descriptor,
and a flag for derived-artifact targets. The model's limitations are thus
recorded explicitly rather than left implicit. Further detail is given in
Supplementary Note~\ref*{supp:sec:supp-ex-inouye}.


% =====================================================================

% =====================================================================
% §5 — Workflow and Evaluation (see paper/sections/05_workflow.tex)
    % =====================================================================
%  §5 — Human–AI Annotation Workflow and Evaluation
% =====================================================================

\section{Human--AI Annotation Workflow and Evaluation}
\label{sec:workflow}

We designed and piloted an annotation workflow in which every annotation is
\emph{auditable}. Each assertion is anchored to the specific text that
supports it; each dimension value is defensible under the schema's
activation conditions; and each judgment that a downstream user might
contest carries an explicit reviewer or AI flag.

The same commitment applies to AI assistance. When an AI system drafts an
annotation, its uncertainties and assumptions are recorded for curator
review rather than silently absorbed into the annotation.
\Cref{sec:workflow-flags} reports the practical effect of that review. The
manual and AI-drafted tracks have different epistemic roles, primary and
secondary knowledge work, respectively, and use distinct flag vocabularies
to record provenance.

We treat the pilot as a \emph{feasibility study}, not as a benchmark.
Six papers and a single annotator do not support claims about
inter-annotator agreement or generalizability. The pilot shows that the
workflow's epistemic commitments can be implemented in practice and that
the schema exposes, rather than hides, its points of strain.

\begin{figure}[tbp]
  \centering
  \begin{tikzpicture}[
      box/.style={draw, rounded corners, align=center, inner sep=3pt,
                  text width=25mm, minimum height=9mm, font=\footnotesize},
      arr/.style={-{Latex[length=2mm]}, semithick}]
    \node[box, fill=black!5] (src) {Source publication (PDF)};
    \node[box, fill=blue!6, below left=12mm and 22mm of src] (m1) {Curator highlights \& callouts};
    \node[box, fill=blue!6, below=6mm of m1] (m2) {Structured JSON record};
    \node[box, fill=blue!6, below=6mm of m2] (m3) {YAML annotation\\(curator ground truth)};
    \node[box, fill=orange!8, below right=12mm and 22mm of src] (a1) {AI draft\\(annotation skill)};
    \node[box, fill=orange!8, below=6mm of a1] (a2) {Curator review\\(review skill)};
    \node[box, fill=orange!8, below=6mm of a2] (a3) {Reviewed YAML annotation};
    \node[box, fill=black!5] (val) at ($(m3)!0.5!(a3) - (0,1.3cm)$) {SHACL validation\\(conditional activation)};
    \draw[arr] (src) -- (m1);
    \draw[arr] (src) -- (a1);
    \draw[arr] (m1) -- (m2);
    \draw[arr] (m2) -- (m3);
    \draw[arr] (a1) -- (a2);
    \draw[arr] (a2) -- (a3);
    \draw[arr] (m3) -- (val);
    \draw[arr] (a3) -- (val);
  \end{tikzpicture}
  \caption{The two annotation tracks. Curator-authored annotations
  (left) form the ground truth. AI-drafted annotations (right) are
  produced under an annotation skill from the source PDF, the schema, and
  the manual annotations as exemplars, then reviewed by a curator under a
  review skill. Both tracks are validated against the same SHACL shapes.
  In every case, the curator, not the model, is the authority for the
  recorded annotation.}
  \label{fig:workflow}
\end{figure}

\subsection{Protocol and Coverage}
\label{sec:workflow-protocol}

Annotations are created through two tracks (\cref{fig:workflow}).

\paragraph{Manual track.}
Four of the six publications, Jossin, Davis, Nelson, and Gupta, were
annotated manually by Vladimir Seplyarskiy, a researcher in human and
population genetics. The curator read each paper, marked relevant
passages in the PDF with highlights and callouts, exported those
annotations with a helper script to a structured JSON record, and mapped
the result, together with a per-paper summary row, to a YAML annotation
conforming to the schema.

When this conversion required normalization, the change was recorded
inline as a \texttt{normalization\_note} with
\texttt{type = ai\_normalization}. For example, the free-text
``localisation and protein expression'' was mapped to the enumeration
values \dimval{Localization} and \dimval{Existence}. This record allows
later audits to compare the original and normalized values. The
manual-track protocol is described in more detail in Supplementary
Note~\ref*{supp:sec:supp-protocol-suite}.

\paragraph{AI-drafted track.}
For two publications, Duerr and Inouye, an AI assistant drafted the
annotations directly from the PDFs under an executable annotation
skill~\cite{anthropic_skills}. The assistant received the schema
(\texttt{dimensions.md}), the ground-truth annotations for the four
manual papers as exemplars, and the source PDFs. It produced YAML
annotations in the same format as the manual annotations.

When the assistant judged a mapping to be uncertain, or concluded that a
judgment call required human review, it recorded the issue inline as
either \texttt{ai\_uncertainty} or \texttt{ai\_assumption}. One of the
authors then reviewed the AI-drafted annotations assertion by assertion.
The reviewer has expertise in genomic data curation and the evidence
model, rather than in genetics domain science. The review assessed source
fidelity and schema conformance, not the underlying science.
The four manual annotations are treated as domain-expert ground truth for
their own papers and as exemplars for AI drafting. The executable skills
and review protocol are described in Supplementary
Notes~\ref*{supp:sec:supp-protocol-suite}
and~\ref*{supp:sec:supp-skills}.

\paragraph{Source anchoring.}
In both tracks, every assertion, that is, the \texttt{assertion} property
of a \GeneticEvidence{} item, carries a \texttt{source\_span}. The span
identifies a specific page and includes a short key phrase, typically fewer
than fifteen words, from the source text. The phrase is long enough for a
reader to locate the passage by search while remaining short enough to
respect fair-use norms for copyrighted material.

Source anchoring is mandatory: an annotation with any assertion lacking a
source span is rejected at validation time. Across the six-paper pilot, all
95 assertions are source-anchored. This is a property of the workflow's
validation requirement, not an observed outcome.

\paragraph{Coverage.}
Dimension coverage across the 28 \GeneticEvidence{} items is summarized in
Supplementary Note~\ref*{supp:sec:supp-coverage-note}
(Table~\ref*{supp:tab:supp-coverage}) and in the repository's coverage
report.\footnote{\url{https://github.com/ForomePlatform/genetic-evidence-model/blob/master/annotations/coverage.md}}

\subsection{Flags and Extension Promotion}
\label{sec:workflow-flags}
\label{sec:workflow-promotion}

The two tracks use small flag vocabularies to make their epistemic
provenance explicit.

\paragraph{Manual track.}
The manual track uses three reviewer flags. A
\texttt{reviewer\_disagreement} records a clear factual correction, such
as the organism-set correction in Davis~2011. A
\texttt{reviewer\_suggestion} records an optional addition whose inclusion
depends on curator preference, such as HPO terms for patient phenotypes in
Davis, Nelson, and Gupta. A \texttt{reviewer\_query} records a case in
which the mapping to the schema is unclear and should not be silently
corrected. For example, the Nelson summary assigned
\texttt{specificity\_of\_phenotype = SNV}, which does not fit the
dimension's semantics.

\paragraph{AI-drafted track.}
Because there is no prior ground truth with which to disagree, the
AI-drafted track uses two flags. An \texttt{ai\_uncertainty} records a
mapping that the assistant judged unreliable and referred for a human
decision, such as a borderline assignment of Mode of Inheritance. An
\texttt{ai\_assumption} records a judgment made without strong textual
support, such as whether cited mouse-model evidence should be represented
as a separate \GeneticEvidence{} item or treated as background context.

Across the six-paper pilot, manual-track annotations carried five queries,
six suggestions, and one disagreement, whereas AI-drafted annotations
carried six uncertainties and six assumptions. Of the nine AI-drafted items
from Duerr and Inouye, none was rejected: one was approved as drafted,
three were approved with non-blocking flags, and five were edited. The
later Inouye draft required fewer edits than the earlier Duerr draft, but
the papers differ in genre and protocol version; this comparison is
therefore suggestive rather than controlled. The edits mainly affected
dimension assignment and credibility calibration rather than the factual
content of assertions. The reviewer also added a small number of
higher-level items, such as a paper's overall thesis or a cited
translational study. Per-paper counts and detailed rationales are given in
Supplementary Note~\ref*{supp:sec:supp-ai-review}.

\paragraph{Extension promotion.}
The workflow adopts a conservative rule for evolving the schema. A
representational gap observed in one paper is recorded only in that
paper's \texttt{candidate\_extensions} block. A gap observed
independently in a second paper is promoted to a first-class dimension or
enumeration value.

In the six-paper pilot, this rule promoted two extensions. \emph{Phenotype
Scale} was first proposed in Jossin and then reused in Davis, Nelson, and
later papers. \emph{Variant Ascertainment} was first proposed in Davis and
then reused in Nelson and Duerr. Both now appear in the vocabulary
(\cref{tab:dimensions-core,tab:dimensions-cond}): Phenotype Scale among
the always-required dimensions and Variant Ascertainment among the
conditional dimensions, where it is required when Target Type is
\dimval{Variant}.

Fifteen other candidates remain unpromoted, pending a second independent
use or an explicit curator decision. The rule also prevents premature
additions. For example, a proposed polarity dimension for negative
assertions in Duerr was withdrawn when review confirmed that the existing
assertion predicate already represents both absence and presence.


% =====================================================================

% =====================================================================
% §6 — Discussion (see paper/sections/06_discussion.tex)
    % =====================================================================
%  §6 — Discussion
%
%  Revised 2026-04-21: subsections collapsed from four to two.
%  §6.1 (positioning) + §6.2 Model Extensions → §6.1 Discussion.
%  §6.3 Limitations + §6.4 Future Work → §6.2 Limitations and
%  Future Work. Labels preserved for all cross-section references
%  (sec:disc-future is referenced from §2).
% =====================================================================

\section{Discussion}
\label{sec:discussion}

Two features of the proposal merit brief emphasis. First, the
always-required vocabulary is small, and conditional activation keeps
the extended vocabulary compact without sacrificing expressive power.
Second, every assertion is anchored to a specific source span in its
source publication, making each annotation independently auditable.

\subsection{Positioning and Model Extensions}
\label{sec:disc-positioning}
\label{sec:disc-extensions}

The released schema formalises the implemented conformance
requirements as SHACL shapes, makes the implemented activation
conditions machine-checkable, and anchors each assertion to its
source text rather than only to a style guide. The flag vocabularies
of \cref{sec:workflow-flags} make provenance explicit: a reader can
tell whether a value was drawn from the author's summary, proposed
by a reviewer, added as an AI assumption, or left deliberately
unpopulated. More broadly, the schema is intended to make literature review tractable as
a \emph{semantic-parsing} problem: the mapping of natural-language claims
to a shared, source-anchored representation. This is not a task a language model can perform in isolation. The target of
each parse is a structured, source-anchored representation, so the quality
of the parse depends on the quality of that target language. We see the present
contribution as one increment along a longer path, building on
decades of biomedical-ontology work (ECO, SEPIO, HPO, IAO, OBI, VRS,
and others), not as a replacement for any of it. A complementary
downstream layer, in which an epistemological type system and
machine-checkable meta-predicates constrain the decision logic that
consumes such evidence, is addressed in~\cite{our_preprint}.

The pilot identified two candidate dimensions that were exercised by
multiple papers and promoted to first-class status, Phenotype Scale and
Variant Ascertainment. It also identified fifteen further candidate
extensions that remain unpromoted under the two-papers rule
(\cref{sec:workflow-promotion}). The pattern of
candidates is summarised in \cref{tab:candidates}; per-candidate
rationale, affected dimensions, and proposed enumeration values
are documented in a companion \texttt{schema/EXTENSIONS.md} in the
repository.\footnote{\url{https://github.com/ForomePlatform/genetic-evidence-model/blob/master/schema/EXTENSIONS.md}}
The Gupta annotation, which identified no candidate extensions, is also a
positive finding. It reports a family with an X-linked splice-site mutation
and careful pedigree segregation, a canonical evidence shape for which the
model was designed. The schema accommodates this evidence without strain.

\begin{table}[tbp]
\centering
\caption{Unpromoted candidate schema extensions surfaced by the
six-paper pilot, grouped by concern. Per-candidate detail is in the
companion \texttt{schema/EXTENSIONS.md}; paper keys are J~=~Jossin,
D~=~Davis, N~=~Nelson, DU~=~Duerr, IN~=~Inouye. Gupta surfaced none.}
\label{tab:candidates}
\small
\begin{tabularx}{\linewidth}{@{}llX@{}}
\toprule
Group & Papers & Count and examples \\
\midrule
Finer structural resolution  & J, N   & 4 (subdomain, interaction, gene-relation, cross-level variant) \\
Polygenic-score machinery    & IN     & 6 (score target, whole-genome aggregate, target composition, epidemiological measurement, cohort sub-block, derived-artefact flag) \\
Effect-size and curation     & DU, D  & 3 (direction of effect, knowledge-domain priority, curator critique) \\
Cross-item and activation scope & DU & 2 (cross-item/translational synthesis, conditional-activation scope) \\
\bottomrule
\end{tabularx}
\end{table}

\subsection{Limitations and Future Work}
\label{sec:disc-limitations}
\label{sec:disc-future}

The pilot has significant limitations. Six papers and a single
annotator do not support claims about inter-annotator agreement or
generalisability; we treat the work as a feasibility study, not a
benchmark. The evaluation is not a clinical
deployment: although the model is motivated by variant
interpretation in clinical-genomics programs, the pilot stops short
of applying the model to a live workflow.

Three near-term directions follow from these limitations. First, annotate a
second corpus selected to exercise the strongest candidate extensions.
Second, generate OWL, SHACL, JSON Schema, and LinkML renderings from a
single source of truth. This would open the schema to the OBO and Bridge2AI
ecosystems and support term-level alignment of the dimension vocabularies
with relevant ontologies, for example by aligning the molecular and cellular
phenotype scales with Gene Ontology biological processes. Third, conduct a
clinical-deployment study to test whether the feasibility claim translates
into utility.

A further limitation revealed by the UMLS crosswalk is representational
rather than empirical. Several GEM tokens are intersective compounds:
\dimval{statistical\_genetics} denotes statistics~$\times$~genetics
(\texttt{C0038215}~$\times$~\texttt{C0017398}), and
\dimval{bioinformatics\_inference} denotes inference~$\times$~bioinformatics.
Relational tokens have the same structure. For example,
\dimval{related\_gene} combines a relation operator with \emph{Genes}
(\texttt{C0017337}). Because the current crosswalk maps each GEM token to
one UMLS concept, it must sometimes use the nearest available anchor rather
than a concept that expresses the full compound meaning. For
\dimval{statistical\_genetics}, that anchor is \emph{Biostatistics},
recorded with the relation \emph{narrower}.

Concept-level post-coordination would address this limitation by mapping a
token to a small set of concepts whose intersection represents its intended
meaning. The UMLS Metathesaurus does not provide such a mechanism at the
aggregate level. As an aggregator, it records the concepts asserted by its
source vocabularies and the links among them, but it does not define a
compositional grammar of its own. Composition remains a feature of individual
source vocabularies, with the compositional grammar of SNOMED~CT being the
best-known example. It cannot be transferred directly across the aggregate:
the same intended intersection may be pre-coordinated in one source, absent
from another, and decomposed differently in a third~\cite{umls2004}.

The problem is particularly clear for \dimval{related\_gene}. UMLS often
contains pre-coordinated specific instances while lacking a generic concept:
gene-relatedness appears in terms such as \emph{modifier},
\emph{homologous genes}, and \emph{linked genes}, whereas the generic sense
of ``related in some way'' is represented only by the bare operator
\emph{Associated with} (\texttt{C0332281}). Thus
\dimval{related\_gene} can be represented as
\texttt{C0332281}~$\times$~\texttt{C0017337}, but not by a single concept.
Extending the crosswalk with an explicit composite relation, represented as
a structured set of concepts, is future work. At present, the intended composition
is preserved in the mapping's curator note; where the relation is the
semantic core, the operator concept itself is recorded as the anchor with
relation \emph{related}.

The crosswalk review also highlighted a representational limitation of the
Resolution axis. Its current enumeration presupposes genomic coordinates, but
genetic evidence may be localized in several coordinate systems. A variant
may be described at the genomic level (HGVS~\texttt{g.}), the transcript
level (\texttt{c.}), or the protein level (\texttt{p.}).

For example, an assertion of the form ``variant $V$ has effect $E$ in
transcript $T$'' is resolved in transcript coordinates. Such assertions are
often tissue-restricted and are characteristic of bioinformatics inference.
The domain mapping underlying the \dimval{PROTEIN\_SUBDOMAIN} candidate is
instead resolved in protein coordinates. In addition, \dimval{VARIANT}
itself has two components: a position and a discrete nucleotide change.

We therefore record \dimval{VARIANT\_IN\_TRANSCRIPT} as a curator-surfaced
candidate value for Resolution. It is separate from the fifteen
pilot-surfaced extensions listed in \cref{tab:candidates} and awaits
independent use in a corpus paper under the promotion rule.

\section{Availability of Code, Data, and Artefacts}
\label{sec:availability}

All artefacts referenced in this paper are open-source and
available in the project repository at
\url{https://github.com/ForomePlatform/genetic-evidence-model}.
This includes the SHACL schema, the human-readable dimension
reference (\texttt{dimensions.md}), per-paper annotations for all
six pilot publications, case reports, the companion
\texttt{schema/EXTENSIONS.md}, the SHACL walkthrough
(\texttt{schema/examples.md}), the annotation and review protocol
suite (\texttt{protocols/}) and the two executable annotation and
review skills (\texttt{skills/}) described in Supplementary Note~SN3,
and the extraction tooling used to
convert PDF highlights and callouts into annotation scaffolds. The
namespace \url{https://w3id.org/genetic-evidence-model/} resolves
to the same repository.

\section*{Acknowledgements}

The idea of a semantic model for the analysis of genetic literature
grew out of 2018 conversations with Prof.\ Shamil Sunyaev; without
his early support this work could not have been completed. The
authors are also grateful to colleagues at the Brigham Genomics
Medicine (BGM) program, where much of the thinking about
clinical-genomics evidence curation took shape, and to
Vladimir Seplyarskiy for his meticulous curation of the four
manually annotated papers in the pilot corpus. This work received
no external funding.


% =====================================================================

% =====================================================================
% §7 — Conclusion (see paper/sections/07_conclusion.tex)
    % =====================================================================
%  §7 — Conclusion
%
%  Target: three to four tight sentences per the
%  section stub in main.tex. Recaps the contribution and lands on
%  trustworthy, AI-ready infrastructure for variant interpretation.
% =====================================================================

\section{Conclusion}
\label{sec:conclusion}

We have introduced a semantic model for genetic evidence from the
biomedical literature: a compact hierarchy of core classes, a
dimensional vocabulary with conditional-activation rules, three of
which are presently machine-enforced in SHACL, and a SHACL encoding
for automated conformance validation. The model is grounded
in an open corpus of six annotated papers and accompanied by a
collaborative workflow that keeps curator-authored reference annotations
distinct from AI-drafted annotations and promotes candidate schema
extensions only when a second independent paper confirms them. The
model is also released with executable annotation and review skills
that let an AI assistant apply it under curator supervision.
Taken together, these elements are offered as a small but concrete
contribution toward trustworthy, AI-ready infrastructure for variant
interpretation: a reference data model and validation schema for
representing genetic evidence from the biomedical literature.


% =====================================================================

% =====================================================================
    \bibliographystyle{plainnat}
    \bibliography{references}
% =====================================================================

\end{document}